% Options for packages loaded elsewhere
\PassOptionsToPackage{unicode}{hyperref}
\PassOptionsToPackage{hyphens}{url}
\PassOptionsToPackage{dvipsnames,svgnames,x11names}{xcolor}
%
\documentclass[
  11pt,
]{article}
\usepackage{amsmath,amssymb}
\usepackage{lmodern}
\usepackage{iftex}
\ifPDFTeX
  \usepackage[T1]{fontenc}
  \usepackage[utf8]{inputenc}
  \usepackage{textcomp} % provide euro and other symbols
\else % if luatex or xetex
  \usepackage{unicode-math}
  \defaultfontfeatures{Scale=MatchLowercase}
  \defaultfontfeatures[\rmfamily]{Ligatures=TeX,Scale=1}
\fi
% Use upquote if available, for straight quotes in verbatim environments
\IfFileExists{upquote.sty}{\usepackage{upquote}}{}
\IfFileExists{microtype.sty}{% use microtype if available
  \usepackage[]{microtype}
  \UseMicrotypeSet[protrusion]{basicmath} % disable protrusion for tt fonts
}{}
\makeatletter
\@ifundefined{KOMAClassName}{% if non-KOMA class
  \IfFileExists{parskip.sty}{%
    \usepackage{parskip}
  }{% else
    \setlength{\parindent}{0pt}
    \setlength{\parskip}{6pt plus 2pt minus 1pt}}
}{% if KOMA class
  \KOMAoptions{parskip=half}}
\makeatother
\usepackage{xcolor}
\IfFileExists{xurl.sty}{\usepackage{xurl}}{} % add URL line breaks if available
\IfFileExists{bookmark.sty}{\usepackage{bookmark}}{\usepackage{hyperref}}
\hypersetup{
  pdftitle={Project Proposal (Week 2), A Short Guide},
  pdfauthor={Course staff},
  colorlinks=true,
  linkcolor={RoyalBlue},
  filecolor={Maroon},
  citecolor={Blue},
  urlcolor={Blue},
  pdfcreator={LaTeX via pandoc}}
\urlstyle{same} % disable monospaced font for URLs
\usepackage[a4paper,margin=2.4cm]{geometry}
\setlength{\emergencystretch}{3em} % prevent overfull lines
\providecommand{\tightlist}{%
  \setlength{\itemsep}{0pt}\setlength{\parskip}{0pt}}
\setcounter{secnumdepth}{-\maxdimen} % remove section numbering
\usepackage{newunicodechar}
\usepackage{textcomp}
\usepackage{amssymb}
\newunicodechar{°}{\textdegree}
\newunicodechar{₂}{\textsubscript{2}}
\newunicodechar{→}{\ensuremath{\rightarrow}}
\newunicodechar{≤}{\ensuremath{\leq}}
\newunicodechar{≥}{\ensuremath{\geq}}
\newunicodechar{≈}{\ensuremath{\approx}}
\newunicodechar{✓}{\ensuremath{\checkmark}}
\newunicodechar{–}{--}
\newunicodechar{—}{---}
\usepackage{etoolbox}
\AtBeginEnvironment{longtable}{\footnotesize}
\AtBeginEnvironment{tabular}{\footnotesize}
\usepackage{fancyhdr}
\pagestyle{fancy}\fancyhf{}
\rhead{\small D7065E\,--\,Project Proposal}
\lhead{\small Week 2 Proposal -- Guide}
\cfoot{\thepage}
\renewcommand{\headrulewidth}{0.4pt}
\usepackage{titlesec}
\titleformat{\section}{\Large\bfseries}{\thesection}{0.6em}{}
\titleformat{\subsection}{\large\bfseries}{\thesubsection}{0.5em}{}
\usepackage{microtype}
\setlength{\emergencystretch}{2em}
\ifLuaTeX
  \usepackage{selnolig}  % disable illegal ligatures
\fi

\title{Project Proposal (Week 2), A Short Guide}
\usepackage{etoolbox}
\makeatletter
\providecommand{\subtitle}[1]{% add subtitle to \maketitle
  \apptocmd{\@title}{\par {\large #1 \par}}{}{}
}
\makeatother
\subtitle{D7065E, what the Week 2 proposal must contain}
\author{Course staff}
\date{}

\begin{document}
\maketitle

\begin{quote}
\textbf{How to use this guide.} It explains what to put in the
\textbf{Week 2 proposal}. See the accompanying worked proposal example.
Keep the proposal to \textbf{1--2 pages}.
\end{quote}

\hypertarget{what-the-proposal-is-and-is-not}{%
\subsection{What the proposal is, and is
not}\label{what-the-proposal-is-and-is-not}}

The proposal exists to agree on a \textbf{direction} and get a
\textbf{green light to start building}. It is \textbf{not} the
architecture: C4 diagrams, requirements tables, interfaces, and tests are
\emph{not} expected yet. Those are built \textbf{iteratively} during
development and end up in the architecture document and final report (see
the report-writing guide and the worked report example). At this stage
all that is needed is the intended build, the reasons for it, and the
open uncertainties. The design \emph{will} change later, that is
expected and good.

\hypertarget{what-it-must-answer}{%
\subsection{What it must answer}\label{what-it-must-answer}}

Six short questions, a paragraph or a few bullets each:

\begin{enumerate}
\def\labelenumi{\arabic{enumi}.}
\tightlist
\item
  \textbf{Use case, and why.} Which use case, and why it was chosen.
\item
  \textbf{Scenario.} The one end-to-end story being aimed for (the
  eventual demo).
\item
  \textbf{What will be built.} A rough sketch of the major pieces: what
  senses, what decides, what acts, \textbf{how data flows through the
  system, where it is stored}, and what users will see. No formal
  diagrams required. Each major piece should run as its \textbf{own
  process}, communicate through defined interfaces, and connect to
  \textbf{BuildSim}.
\item
  \textbf{The autonomous part.} Roughly what decides what, and the
  approach to \emph{start} with. Rule-based, optimisation, ML, or
  LLM are all fine, \emph{sophistication earns no marks}, so start
  simple.
\item
  \textbf{Biggest unknowns \& risks.} What is uncertain or might not
  work. Honesty here is the point.
\item
  \textbf{Physical simulation approach.} How environmental conditions
  and occupancy will be generated, what physical model, dataset, or
  hybrid approach will drive the sensor values. This makes the
  simulator a first-class part of the project, not an implementation
  detail; the next section gives examples.
\item
  \textbf{First step \& rough plan.} The riskiest thing to prototype
  first to test the idea, plus a rough sense of the weeks ahead.
\end{enumerate}

Include, at the top of the proposal, a link to the team's
\textbf{GitHub repository} where the project will be developed.

\hypertarget{simulating-the-building}{%
\subsection{Simulating the building}\label{simulating-the-building}}

BuildSim provides the rooms but \textbf{not their physics}, so the
proposal should say how the \textbf{physical processes} that produce
sensor values and the \textbf{movement of people} that drives them are
simulated.

\begin{itemize}
\tightlist
\item
  \textbf{Physical processes}, sketch how each sensor value will be
  generated: a simple physics-based model (for example a heat balance
  for temperature, or a mass balance for CO₂ that rises with occupancy
  and falls with ventilation), replay of a public dataset, or a hybrid.
  It need not be exact, it must have realistic trends and respond to
  the actuators so the control loop closes.
\item
  \textbf{Movement of people (occupancy)}, occupancy usually drives
  everything else (heat, CO₂, lighting). Describe the pattern to
  model: a weekday schedule (arrivals, lunch, departures), meeting-room
  bookings, empty weekends, and some randomness so no two days are
  identical.
\end{itemize}

\textbf{Examples of approaches.} Simple physics, a heat balance for
temperature, or CO₂ rising about 40 ppm per person per hour and decaying
with ventilation. Data replay, feed a public dataset through the
sensors (e.g.~the UCI \emph{Occupancy Detection} set, or \emph{ASHRAE}
building-energy data). Hybrid, replay real data with an injected
event (e.g.~a real temperature trace plus a simulated fire). For
occupancy, a schedule such as arrive 07:00--09:00, lunch 11:30--13:00,
leave 16:00--18:00, with meeting bookings and weekend gaps.

No equations are needed yet, just the intended approach.

\hypertarget{how-it-is-approved}{%
\subsection{How it is approved}\label{how-it-is-approved}}

The proposal is approvable when the \textbf{direction is clear and
feasible}, not when the design is complete. Approval is a green light
to start, not a grade. What follows is the architecture document, which
grows alongside the code; a brief ``changes since the proposal'' note
in the final report, explaining what was revised and why, is something
the examiner looks for.

\end{document}
